\quad O Selection Sort é um algoritmo de ordenação baseado na seleção
sucessiva do menor elemento do vetor. Seu funcionamento consiste em
dividir mentalmente o vetor em duas partes: uma parte inicialmente
desordenada e uma parte que já foi ordenada. A cada iteração, o algoritmo
procura o menor elemento da parte desordenada e o coloca na posição
correspondente da parte ordenada.

\quad Inicialmente, considera-se que nenhum elemento está ordenado.
Durante cada iteração, o algoritmo percorre os elementos restantes do
vetor e identifica aquele que possui o menor valor. Após encontrar esse
elemento, ele é trocado com o primeiro elemento da parte ainda
desordenada. Dessa forma, a parte ordenada cresce progressivamente até
que todos os elementos estejam em suas posições corretas.

\begin{lstlisting}[caption={Implementação do Selection Sort}]
void selectionSort(int* vetor, int tamanho) {
    for (int i = 0; i < tamanho - 1; i++) {
        int menor = i; 

        for (int j = i + 1; j < tamanho; j++) {
            if (vetor[j] < vetor[menor]) {
                menor = j; 
            }
        }

        int temp = vetor[i];
        vetor[i] = vetor[menor];
        vetor[menor] = temp;
    }
}
\end{lstlisting}

\quad No primeiro laço, a variável $i$ representa o início da parte ainda
não ordenada do vetor. A variável \texttt{menor} armazena a posição do
menor elemento encontrado durante a busca.

\quad O segundo laço percorre os elementos que ainda não foram ordenados,
comparando seus valores com o elemento armazenado na posição
\texttt{menor}. Quando um elemento menor é encontrado, sua posição passa
a ser armazenada nessa variável.

\quad Ao final de cada iteração, o menor elemento encontrado é trocado
com o elemento que ocupa a posição $i$. Assim, esse elemento passa a
fazer parte da região ordenada do vetor. O processo é repetido até que
restem apenas dois elementos ou até que todo o vetor esteja ordenado.
